\section{The Class Equation and Center of \texorpdfstring{$p$-Groups}{p-Groups}}

\subsection{Orbit-Stabilizer Applied to Conjugation}

By the Orbit-Stabilizer Theorem, for any $a \in G$, there exists a bijection between the conjugacy class $\operatorname{cl}(a)$ and the set of cosets $G/N(a)$, for all $a \in G$. Consequently:
\[
G = \bigcup_{a \in G} \operatorname{cl}(a).
\]
If $G$ is a finite group, then:
\begin{align*}
|\operatorname{cl}(a)| &= \left|\frac{G}{N(a)}\right| = \frac{|G|}{|N(a)|}, \\
|G| &= \sum_{a \in G} |\operatorname{cl}(a)|.
\end{align*}

\begin{proposition}
For any element $a \in G$:
\[
\operatorname{cl}(a) = \{a\} \iff a \in Z(G),
\]
where $Z(G) = \{z \in G \mid zg = gz \; \forall g \in G\}$ is the center of $G$.
\end{proposition}

\begin{proof}
By definition of the conjugacy class:
\begin{align*}
\operatorname{cl}(a) = \{a\} &\iff gag^{-1} = a \quad \forall g \in G \\
&\iff ga = ag \quad \forall g \in G \\
&\iff a \in Z(G).
\end{align*}
\end{proof}

\subsection{The Class Equation}

Separating the central elements (for which $|\operatorname{cl}(a)| = 1$) from the non-central elements yields the \textbf{Class Equation}:
\[
|G| = |Z(G)| + \sum_{a \notin Z(G)} \left|\frac{G}{N(a)}\right| \longrightarrow (1)
\]

\begin{theorem}
If $|G| = p^n$, where $p$ is a prime and $n \ge 1$, then:
\[
|Z(G)| \ge p.
\]
In other words, every non-trivial finite $p$-group has a non-trivial center.
\end{theorem}

\begin{proof}
Let $|Z(G)| = r$. By the class equation (1):
\[
|G| = |Z(G)| + \sum_{a \notin Z(G)} \left|\frac{G}{N(a)}\right|.
\]
For any $a \notin Z(G)$, $N(a) \neq G$, meaning $N(a)$ is a proper subgroup of $G$. By Lagrange's Theorem, $|N(a)|$ divides $|G| = p^n$, so:
\[
|N(a)| = p^{\alpha_a} \quad \text{with } \alpha_a < n.
\]
Then:
\[
\left|\frac{G}{N(a)}\right| = \frac{|G|}{|N(a)|} = \frac{p^n}{p^{\alpha_a}} = p^{n - \alpha_a}, \quad \text{where } n - \alpha_a > 0.
\]
This implies:
\[
p \mid p^{n - \alpha_a} \implies p \mid \left|\frac{G}{N(a)}\right| \implies p \mid \sum_{a \notin Z(G)} \left|\frac{G}{N(a)}\right|.
\]
Since $|G| = p^n$, we also have $p \mid |G|$. Rewriting the class equation:
\[
|Z(G)| = |G| - \sum_{a \notin Z(G)} \left|\frac{G}{N(a)}\right|.
\]
Since $p$ divides both terms on the right-hand side:
\[
p \mid |Z(G)| \implies p \mid r.
\]
Since the identity element $e \in Z(G)$, the center is never empty, so $|Z(G)| \ge 1$. Because $p$ divides $r$, we must have:
\[
r \ge p \implies |Z(G)| \ge p.
\]
\end{proof}

\begin{theorem}
A group of order $p^2$, where $p$ is prime, is abelian.
\end{theorem}

\begin{proof}
Let $|G| = p^2$. By the previous theorem:
\[
|Z(G)| \ge p.
\]
By Lagrange's Theorem, $|Z(G)|$ must divide $|G| = p^2$. Therefore, the only possibilities are:
\[
|Z(G)| = p \quad \text{or} \quad |Z(G)| = p^2.
\]
If $|Z(G)| = p$, then:
\[
\frac{|G|}{|Z(G)|} = \frac{p^2}{p} = p.
\]
Since any group of prime order is cyclic, the quotient group $G/Z(G)$ is cyclic. But it is a standard result that:
\[
G/Z(G) \text{ is cyclic} \implies G \text{ is abelian}.
\]
If $G$ is abelian, every element commutes with all other elements, so:
\[
G = Z(G) \implies |G| = |Z(G)| = p^2,
\]
which contradicts the assumption that $|Z(G)| = p$.

Therefore:
\[
|Z(G)| \neq p \quad \text{(contradiction)}.
\]
So we must have:
\[
|Z(G)| = p^2 = |G| \implies G = Z(G).
\]
Hence, $G$ is abelian.
\end{proof}

\begin{remark}
If $H \le G$ such that $H \subseteq Z(G)$, then $H \trianglelefteq G$ (any subgroup contained in the center is normal in $G$).
\end{remark}
